\documentclass[11pt,a4paper]{article}
\usepackage[margin=20mm]{geometry}
\usepackage[T1]{fontenc}
\usepackage[utf8]{inputenc}
\usepackage{xcolor}
\usepackage{printsudoku}
\usepackage{solvesudoku}

\definecolor{puzzleblue}{HTML}{1E4E79}
\title{Sudoku Puzzle and Solution}
\author{Your Name}
\date{\today}

\begin{document}
\maketitle

\noindent\textcolor{puzzleblue}{\bfseries Puzzle file:} `puzzle.sud'\par
\medskip

\begin{center}
  {\large\bfseries Solve it yourself}\par\medskip
  \cluefont{\Large}\cellsize{2.2\baselineskip}
  \sudoku{puzzle.sud}
\end{center}

\newpage
\section*{Solver output}
\cluefont{\small}\cellsize{1.2\baselineskip}
\sudokusolve{puzzle.sud}

\section*{Workflow}
Keep one puzzle per `.sud` file.  The first nine lines contain clues, with a
dot for each blank cell; comments after those lines are ignored by
`printsudoku`.  Use the solver view as a checking aid while preparing a
worksheet or a puzzle collection.
\end{document}
